% Average comparisons for a successful search. Both axes are logarithmic,
% because the four series do not share an order of growth: three of them are
% logarithmic and one is linear.
%
% The separation is the whole argument for balancing. On sorted keys the plain
% tree needs hundreds of times more comparisons than the balanced one, and the
% factor grows with every key added.
\begin{tikzpicture}
  \begin{axis}[
      avlaxis,
      xmode        = log,
      ymode        = log,
      log basis x  = 10,
      xlabel       = {number of keys $n$},
      ylabel       = {mean comparisons per search},
      xmin         = 80, xmax = 1.3e4,
      ymin         = 4, ymax = 1.2e4,
      % Below the axis: every corner inside the frame has a curve running
      % through it, so an inset legend would sit on top of the data.
      legend style = {at = {(0.5, -0.30)}, anchor = north, draw = rfaint},
      legend columns = 2,
      legend cell align = left,
    ]

    \addplot[mark = *, mark size = 1.5pt, thick, color = rmid, dashed]
      table[x = n, y = bstsorted] {data/search_cost.dat};
    \addlegendentry{plain tree, sorted keys}

    \addplot[avlsecond] table[x = n, y = bstrand] {data/search_cost.dat};
    \addlegendentry{plain tree, random keys}

    \addplot[avlmeasured] table[x = n, y = avlrand] {data/search_cost.dat};
    \addlegendentry{AVL tree, random keys}

    \addplot[no marks, thick, dotted, color = rkink]
      table[x = n, y = avlsorted] {data/search_cost.dat};
    \addlegendentry{AVL tree, sorted keys}

  \end{axis}
\end{tikzpicture}
